\chapter{Jeux sous forme extensive}

Les jeux sous forme normale donne une représentation adéquate de
joueurs effectuant des choix de stratégies simultanés. Dans ces
jeux, aucun jouer n'a d'information sur les stratégies utilisées
par les autres au moment de faire ses choix. Dans de nombreux cas
concrets, les joueurs effectuent leurs choix en fonction d'une
situation observée qui dépend en particulier des choix précédents
des autres joueurs. Par exemple c'est le cas pour une enchère
croissante, dans une négociation, ou dans le jeu d'échecs...

Pour représenter ces possibilités, nous allons généraliser la
représentation sous forme d'arbres et les  ensembles d'information
vus au chapitre de la décision dynamique. La seule différence est
que nous avons plusieurs joueurs là où nous n'avions qu'un seul
agent.

\section{Exemples préliminaires}
Nous traitons une série d'exemples avant de définir la notion de
jeu sous forme extensive.

\subsection{Jeu d'entrée} Une firme $M$ est en situation de
monopole sur un marché. Une autre firme $E$ peut décider d'entrer
($e$) ou non ($n$) sur le marché. Si la firme $E$ décide d'entrer,
la firme $M$ peut alors soit combattre ($c$), soit s'accommoder
($a$). Les paiements pour $M$ et pour $A$ sont $(2,0)$ si $E$
n'entre pas, $(-1,-1)$ si $E$ entre et $M$ combat, et $(1,1)$ si
$E$ entre et $M$ s'accommode. On représente ce jeu par l'arbre suivant:\\[1cm]
\begin{tabular}{c@{\hspace{1cm}}c@{\hspace{1cm}}c@{\hspace{1cm}}c@{\hspace{1cm}}}
&&\circlenode{a}{\makebox[0.2cm][c]{\small $E$}}&\\[0.7cm]
&\circlenode{b}{\makebox[0.2cm][c]{\small $M$}}&&\rnode{c}{\small$2$}\\
&&&{\small$0$}\\[0.4cm]
\rnode{d}{\small$-1$}&&\rnode{e}{\small$1$}&\\
{\small$-1$}&&{\small$1$}&\\[0.7cm]
\ncline[linewidth=1pt,linecolor=yellow]{-}{a}{b}\bput[1pt]{0}(0.5){\small$e$}%
\ncline[linewidth=1pt,linecolor=yellow]{-}{a}{c}\aput[1pt]{0}(0.5){\small$n$}%
\ncline[linewidth=1pt,linecolor=yellow]{-}{b}{d}\bput[1pt]{0}(0.5){\small$c$}%
\ncline[linewidth=1pt,linecolor=yellow]{-}{b}{e}\aput[1pt]{0}(0.5){\small$a$}%
%\rput{45}{\psellipse[linecolor=green](4.2,1)(1.3,0.3)}
\end{tabular}

\subsection{Compétition de Stackelberg}
Supposons qu'une firme, 1, ait la possibilité de choisir son
niveau de production $q_1$ avant que la firme $2$ ne fixe le sien
$q_2$. Pour fixer les idées, supposons que $q_i \in \{0,1,2\}$,
que la fonction de demande est $p=\max\{3-q_1+q_2,0\}$, et que le
co\^ut de production est nul pour chaque firme.
On représente ce jeu par l'arbre suivant:\\[0.8cm]
\begin{center}
\begin{tabular}{c@{\hspace{1cm}}c@{\hspace{1cm}}c@{\hspace{1cm}}c@{\hspace{1cm}}c@{\hspace{1cm}}c@{\hspace{1cm}}%
c@{\hspace{1cm}}c@{\hspace{1cm}}c@{\hspace{1cm}}}
&&&&\circlenode{a}{\makebox[0.2cm][c]{\small $1$}}&&&&\\[1cm]
&\circlenode{b}{\makebox[0.2cm][c]{\small $2$}}&&&
\circlenode{c}{\makebox[0.2cm][c]{\small $2$}}&&&
\circlenode{d}{\makebox[0.2cm][c]{\small $2$}}&\\[1cm]
\rnode{e}{\small$0$}&\rnode{f}{\small$0$}&\rnode{g}{\small$0$}&%
\rnode{h}{\small$2$}&\rnode{i}{\small$1$}&\rnode{j}{\small$0$}&%
\rnode{k}{\small$2$}&\rnode{l}{\small$0$}&\rnode{m}{\small$0$}\\
\small$0$&\small$2$&\small$2$&%
\small$0$&\small$1$&\small$0$&%
\small$0$&\small$0$&\small$0$\\%
\ncline[linewidth=1pt,linecolor=yellow]{-}{a}{b}\bput[1pt]{0}(0.5){\small$0$}%
\ncline[linewidth=1pt,linecolor=yellow]{-}{a}{c}\aput[1pt]{0}(0.5){\small$1$}%
\ncline[linewidth=1pt,linecolor=yellow]{-}{a}{d}\aput[1pt]{0}(0.5){\small$2$}%
\ncline[linewidth=1pt,linecolor=yellow]{-}{b}{e}\bput[1pt]{0}(0.5){\small$0$}%
\ncline[linewidth=1pt,linecolor=yellow]{-}{b}{f}\bput[1pt]{0}(0.5){\small$1$}%
\ncline[linewidth=1pt,linecolor=yellow]{-}{b}{g}\aput[1pt]{0}(0.5){\small$2$}%
\ncline[linewidth=1pt,linecolor=yellow]{-}{c}{h}\bput[1pt]{0}(0.5){\small$0$}%
\ncline[linewidth=1pt,linecolor=yellow]{-}{c}{i}\bput[1pt]{0}(0.5){\small$1$}%
\ncline[linewidth=1pt,linecolor=yellow]{-}{c}{j}\aput[1pt]{0}(0.5){\small$2$}%
\ncline[linewidth=1pt,linecolor=yellow]{-}{d}{k}\bput[1pt]{0}(0.5){\small$0$}%
\ncline[linewidth=1pt,linecolor=yellow]{-}{d}{l}\bput[1pt]{0}(0.5){\small$1$}%
\ncline[linewidth=1pt,linecolor=yellow]{-}{d}{m}\aput[1pt]{0}(0.5){\small$2$}%
\end{tabular}
\end{center}

\subsection{Pile ou face}

Le joueur 1 choisit pile ou face, puis le joueur 2 choisit pile ou
face. Le joueur 2 gagne si les choix sont égaux, et perd si les
choix sont différents (c'est un jeu à somme nulle).
On représente ce jeu par l'arbre suivant:\\[0.5cm]
\begin{center}
\begin{tabular}{c@{\hspace{0.5cm}}c@{\hspace{0.5cm}}c@{\hspace{0.5cm}}%
c@{\hspace{0.5cm}}c@{\hspace{0.5cm}}c@{\hspace{0.5cm}}%
c@{\hspace{0.5cm}}}
&&&\circlenode{a}{\makebox[0.2cm][c]{\small $1$}}&&&\\[1cm]
&\circlenode{b}{\makebox[0.2cm][c]{\small $2$}}&&&&
\circlenode{c}{\makebox[0.2cm][c]{\small $2$}}&\\[1cm]
\rnode{d}{\small$1$}&&\rnode{e}{\small$-1$}&&\rnode{f}{\small$-1$}%
&&\rnode{g}{\small$1$}\\
\small$-1$&&\small$1$&&\small$1$&&\small$-1$\\%
\ncline[linewidth=1pt,linecolor=yellow]{-}{a}{b}\bput[1pt]{0}(0.5){\small$P$}%
\ncline[linewidth=1pt,linecolor=yellow]{-}{a}{c}\aput[1pt]{0}(0.5){\small$F$}%
\ncline[linewidth=1pt,linecolor=yellow]{-}{b}{d}\bput[1pt]{0}(0.5){\small$P$}%
\ncline[linewidth=1pt,linecolor=yellow]{-}{b}{e}\aput[1pt]{0}(0.5){\small$F$}%
\ncline[linewidth=1pt,linecolor=yellow]{-}{c}{f}\bput[1pt]{0}(0.5){\small$P$}%
\ncline[linewidth=1pt,linecolor=yellow]{-}{c}{g}\aput[1pt]{0}(0.5){\small$F$}%
\end{tabular}
\end{center}

\subsection{Pile ou Face sans observation}
On suppose maintenant que 2 ne connaît pas le choix de 1 lorsqu'il
effectue son choix. Ceci peut être représenté à l'aide d'un
ensemble d'information :
\begin{center}
\begin{tabular}{c@{\hspace{0.5cm}}c@{\hspace{0.5cm}}c@{\hspace{0.5cm}}%
c@{\hspace{0.5cm}}c@{\hspace{0.5cm}}c@{\hspace{0.5cm}}%
c@{\hspace{0.5cm}}}
&&&\circlenode{a}{\makebox[0.2cm][c]{\small $1$}}&&&\\[1cm]
&\circlenode{b}{\makebox[0.2cm][c]{\small $2$}}&&&&
\circlenode{c}{\makebox[0.2cm][c]{\small $2$}}&\\[1cm]
\rnode{d}{\small$1$}&&\rnode{e}{\small$-1$}&&\rnode{f}{\small$-1$}%
&&\rnode{g}{\small$1$}\\
\small$-1$&&\small$1$&&\small$1$&&\small$-1$\\%
\ncline[linewidth=1pt,linecolor=yellow]{-}{a}{b}\bput[1pt]{0}(0.5){\small$P$}%
\ncline[linewidth=1pt,linecolor=yellow]{-}{a}{c}\aput[1pt]{0}(0.5){\small$F$}%
\ncline[linewidth=1pt,linecolor=yellow]{-}{b}{d}\bput[1pt]{0}(0.5){\small$P$}%
\ncline[linewidth=1pt,linecolor=yellow]{-}{b}{e}\aput[1pt]{0}(0.5){\small$F$}%
\ncline[linewidth=1pt,linecolor=yellow]{-}{c}{f}\bput[1pt]{0}(0.5){\small$P$}%
\ncline[linewidth=1pt,linecolor=yellow]{-}{c}{g}\aput[1pt]{0}(0.5){\small$F$}%
\ncline[linewidth=2pt,linecolor=green,linestyle=dotted]{-}{b}{c}% ensemble d'information
\end{tabular}
\end{center}
L'ensemble d'information du joueur 2 comprend ses deux n{\oe}uds
de décision. Ceci signifie que le joueur 2 ne différencie pas ces
deux n{\oe}uds. Il doit effectuer un choix indépendant de celui
effectué précédemment par le joueur 1.
\subsection{Pile ou Face contre la nature} Cette fois ci, il
n'y a pas de joueur 2, mais le joueur 1 joue contre la Nature qui
tire une pièce de monnaie, Pile ou Face, avec probabilités $\lambda$ et $1-\lambda$.\\[0.5cm]
\begin{center}
\begin{tabular}{c@{\hspace{0.5cm}}c@{\hspace{0.5cm}}c@{\hspace{0.5cm}}%
c@{\hspace{0.5cm}}c@{\hspace{0.5cm}}c@{\hspace{0.5cm}}%
c@{\hspace{0.5cm}}}
&&&\circlenode{a}{\makebox[0.2cm][c]{\small $\red N$}}&&&\\[1cm]
&\circlenode{b}{\makebox[0.2cm][c]{\small $1$}}&&&&
\circlenode{c}{\makebox[0.2cm][c]{\small $1$}}&\\[1cm]
\rnode{d}{\small$1$}&&\rnode{e}{\small$0$}&&\rnode{f}{\small$0$}%
&&\rnode{g}{\small$1$}\\
\ncline[linewidth=1pt,linecolor=yellow]{-}{a}{b}\bput[1pt]{0}(0.6){\small$P$}%
\ncline[linewidth=1pt,linecolor=yellow]{-}{a}{c}\aput[1pt]{0}(0.6){\small$F$}%
\ncline[linewidth=1pt,linecolor=yellow]{-}{a}{b}\bput[1pt]{0}(0.3){\small\red$\lambda$}%
\ncline[linewidth=1pt,linecolor=yellow]{-}{a}{c}\aput[1pt]{0}(0.3){\small\red$1-\lambda$}%
\ncline[linewidth=1pt,linecolor=yellow]{-}{b}{d}\bput[1pt]{0}(0.5){\small$P$}%
\ncline[linewidth=1pt,linecolor=yellow]{-}{b}{e}\aput[1pt]{0}(0.5){\small$F$}%
\ncline[linewidth=1pt,linecolor=yellow]{-}{c}{f}\bput[1pt]{0}(0.5){\small$P$}%
\ncline[linewidth=1pt,linecolor=yellow]{-}{c}{g}\aput[1pt]{0}(0.5){\small$F$}%
\ncline[linewidth=2pt,linecolor=green,linestyle=dotted]{-}{b}{c}% ensemble d'information
\end{tabular}
\end{center}
$N$ représente la nature. On a représenté les probabilités
$\lambda, 1-\lambda$ des choix pouvant être faits par la nature.

\section{Jeu sous forme extensive $\Gamma$}

Voici la définition formelle d'un jeu sous forme extensive. On ne
demandera pas de la connaître par c{\oe}ur. Cependant, il est
important de savoir reconnaître un tel jeu et de savoir
représenter une situation stratégique par un jeu sous forme
extensive.
\begin{definition}
Un jeu sous forme extensive $\Gamma$ est donné par :

\begin{itemize}
\item $I$ ensemble fini de joueurs,  $N\not\in I$ (nature);
\item $T$ ensemble fini de n{\oe}uds, $Z\subset T$ n{\oe}uds terminaux, et pour
$t\not\in Z$ :
\begin{itemize}
\item Le joueur ou la nature $i(t)\in \{N\}\cup I$ qui joue en $t$ ;
\item L'ensemble de choix possibles $A(t)$ ;
\item Le n{\oe}ud successeur résultant du choix $a$, $N(t,a)$ ;
\item Une distribution $D(t)$ sur $A(t)$ si $i(t)= N$.
\end{itemize}
\item $u_i:Z\to\reals$ fonctions d'utilité de $i$ ;
\item $h(t)$ classe d'équivalence de $t$ pour l'information :
$$
x'\in h(x) \implies i(x')= i(x),~A(x')=A(x), h(x')=h(x)
$$
\end{itemize}

Afin que le jeu ait une structure d'arbre, on requiert que les
éléments vérifient les propriétés suivantes.

\begin{itemize}
\item $s(t)=\{N(t,a),a\in A(t)\}$ ensemble des successeurs directs de $t$. ($\emptyset$ si $t\in Z$).
\item $S(t)= s(t)\cup s(s(t)) \cup s(s(s(t)))...$ tous les successeurs de $t$.
\end{itemize}
\begin{itemize}
\item Il n'y a pas de cycle : $\forall t, t\not\in S(t)$
\item N{\oe}ud initial $\exists\tilde t, S(\tilde t)\cup\{\tilde t\}=T$
\item Chaque n{\oe}ud non initial a un unique prédécesseur:
$\forall t\neq \tilde t,~\exists ! t',~t\in s(t')$
\end{itemize}
\end{definition}
En exercice, on pourra définir les composants $T$, $Z$, etc qui
représentent les jeux précédents, puis vérifier que ces jeux ont
bien les propriétés de structure d'arbre.

\subsection{Hypothèse de mémoire parfaite}

Il est naturel de supposer qu'un joueur ne peut oublier de
l'information au cours du déroulement du jeu. Notamment, ce joueur
doit se rappeler tous ses choix précédents. C'est l'hypothèse dite
de mémoire parfaite.

\begin{definition}
Le jeu $\Gamma$ est à mémoire parfaite si pour tout $x, y, z$ tels
que $i(x)=i(y)=i(z)$, $h(y)=h(z)$ et $a$ tels que $z\in
S(N(x,a))\cup\{N(x,a)\}$, il existe $w$ tel que :
\begin{itemize}
\item $h(w)=h(x)$ ;
\item $y\in S(N(w,a))\cup\{N(w,a)\}$.
\end{itemize}
\end{definition}

A titre d'exemple, considérons le cas du conducteur distrait. Une
autoroute comprend deux sorties. Afin de se rendre chez lui, un
automobiliste doit prendre la seconde sortie. Cependant, on
suppose que le conducteur est distrait, c'est-à-dire qu'au moment
de décider de sortir de l'autoroute ou d'y rester, il est
incapable de déterminer s'il se trouve à la première ou a la
deuxième sortie. Sauriez-vous représenter cette situation par un
arbre de jeu ? L'hypothèse de mémoire parfaite est-elle vérifiée ?
Pourquoi ?

Sauf exception, on ne considérera des jeux avec mémoire parfaite. %

\subsection{Stratégies pures}
Comme dans le cas de la décision dynamique, une stratégie pour un
joueur est une fonction de ces ensembles d'information vers ses
choix.

\begin{definition}
 Une stratégie pure pour $i\in I$ dans $\Gamma$ est une
application $s_i\colon \{t,i(t)=i\}\to\cup_{t}A(t)$ telle que:
\begin{itemize}
\item $\forall t, s_i(t)\in A(t)$ ;
\item $s_i(t) = s_i(t')$ si $t\in h(t')$.
\end{itemize}
\end{definition}

Remarquons que cette définition permet facilement d'énumérer les
stratégies du joueur. Cependant, deux stratégies peuvent toujours
donner le même résultat. On dira alors qu'elles sont équivalentes
(nous reviendrons plus tard sur cette notion).

Comme exercice, et à titre de rappel des problèmes de décision
dynamiques, quelles sont les stratégies du joueur 1 dans le jeu
suivant ? Quelles stratégies sont équivalentes ?

\begin{tabular}{c@{\hspace{1cm}}c@{\hspace{1cm}}c@{\hspace{1cm}}c@{\hspace{1cm}}}
&&\circlenode{a}{\makebox[0.2cm][c]{\small $1$}}&\\[0.7cm]
&\circlenode{b}{\makebox[0.2cm][c]{\small $1$}}&&\rnode{c}{\small$2$}\\
&&&{}\\[0.4cm]
\rnode{d}{\small$-1$}&&\rnode{e}{\small$1$}&\\
{}&&{}&\\[0.7cm]
\ncline[linewidth=1pt,linecolor=yellow]{-}{a}{b}\bput[1pt]{0}(0.5){\small$a$}%
\ncline[linewidth=1pt,linecolor=yellow]{-}{a}{c}\aput[1pt]{0}(0.5){\small$b$}%
\ncline[linewidth=1pt,linecolor=yellow]{-}{b}{d}\bput[1pt]{0}(0.5){\small$c$}%
\ncline[linewidth=1pt,linecolor=yellow]{-}{b}{e}\aput[1pt]{0}(0.5){\small$d$}%
%\rput{45}{\psellipse[linecolor=green](4.2,1)(1.3,0.3)}
\end{tabular}

\subsection{Stratégies aléatoires}

Comme dans le cas des jeux sous forme normale, nous avons besoin
de représenter les possibilités de jouer de manière aléatoire.
Deux types de stratégies représentent des choix aléatoires des
joueurs.
\begin{definition}
Une stratégie mixte est une loi de probabilité sur les stratégies
pures.
\end{definition}


\begin{definition}
Une stratégie de comportement est une application $\sigma_i\colon
\{t,i(t)=i\}\to\cup_{t}\Delta(A(t))$ telle que :
\begin{itemize}
\item $\forall t, \sigma_i(t)(a)>0\implies a\in A(t)$ ;
\item $\sigma_i(t) = \sigma_i(t')$ si $t\in h(t')$.
\end{itemize}
\end{definition}

Dans une stratégie mixte, le joueur choisit aléatoirement ``au
début du jeu'' la stratégie pure qu'il utilisera par la suite. Une
fois qu'il a choisi cette stratégie, il joue en suivant cette
règle déterministe de décisions. Une stratégie de comportements,
spécifie un choix aléatoire à chaque n{\oe}ud de décision, on donc
un aléa a chaque fois qu'il s'agit de prendre une décision. L'aléa
porte sur la prochaine action choisie, et non pas sur une règle
globale de comportement.

\subsection{\'Equivalence des stratégies mixtes et de comportement}

Deux stratégies $\sigma_i$ et $\sigma_i'$ du joueur $i$ (pures,
mixtes ou de comportement) sont équivalentes lorsque pour toutes
stratégies $\sigma_{-i}$ des autres joueurs (idem),
$(\sigma_i,\sigma_{-i})$ et $(\sigma'_i,\sigma_{-i})$ induisent la
m\^eme distribution sur les n{\oe}uds terminaux.

Comme nous l'avons vu, des stratégies pures peuvent en particulier
être équivalentes entre elles.

\subsection{Théorème de Kuhn}

\begin{theoreme}[de Kuhn]
Dans tout jeu à mémoire parfaite,
\begin{enumerate}
\item Toute stratégie mixte admet une stratégie de comportement
qui lui est équivalente ;
\item Toute stratégie de comportement
admet une stratégie mixte qui lui est équivalente.
\end{enumerate}
\end{theoreme}

Les deux manières de modéliser la capacité de jouer aléatoirement
-- en utilisant les stratégies mixtes ou les stratégies de
comportement -- sont donc équivalentes dans les jeux à mémoire
parfaite. On utilisera en pratique la représentation la plus
commode selon les cas.%

Ce résultat n'est vrai que pour les jeux à mémoire parfaite.
Reprenez l'exemple du conducteur distrait. Toutes les stratégies
mixtes sont-elles équivalentes à des stratégies de comportement ?
Le contraire est-il vrai ?

\section{Forme normale et équilibres de Nash}
Pour ``résoudre'' des jeux en forme extensive, nous allons
utiliser le concept d'équilibre de Nash de la forme normale, puis
montrer que la forme extensive permet de faire des prédictions
plus fines (équilibre parfait dans les sous-jeux).

\subsection{Forme normale et forme normale réduite}
La forme normale associée à
$\Gamma$ est le jeu sous forme normale $G = (I,(S_i)_i,(g_i)_i)$,
avec
$$
g_i(s) = \E_{P_s} u_i(z)
$$
où $P_s$ est la probabilité sur les n{\oe}uds terminaux induite
par le profil de stratégies (pures) $s$.

Soit $S^r_i$ l'ensemble des classes d'équivalences de stratégies
du joueur $i$. Un élément de $S^r_i$ est appelé
une stratégie réduite du joueur $i$.\\
La forme normale réduite de $\Gamma$ est le jeu sous forme normale
$G^r = (I,(S^r_i)_i,(g^r_i))$ avec
$$
g^r_i (s^r)= g_i(s)
$$
où $s^r = (s^r_j)_j$, $s = (s_j)_j$, $s_j\in s^r_j$ pour tout $j$.
(Le choix de $s_j$ dans $s_j^r$ est indifférent.)

\subsection{\'Equilibre de Nash}

\begin{definition}
Un équilibre de Nash d'un jeu sous forme extensive est un
équilibre de Nash de la forme normale associée.
\end{definition}

{\bf Remarque :} Tout les équilibres de Nash d'un jeu sont obtenus
dans la forme normale réduite.

{\bf Remarque :} Lorsqu'on recherche les équilibres de Nash d'un
jeu sous forme extensive, on utilise la notion de stratégies
mixtes. Cette notion est ici plus utile que celle de stratégies de
comportement.

{\bf Exemple :}
Quels sont les équilibres de Nash du jeu d'entrée (en stratégies pures ou mixtes)~?\\
\begin{center}
\begin{tabular}{c@{\hspace{1cm}}c@{\hspace{1cm}}c@{\hspace{1cm}}c@{\hspace{1cm}}}
&&\circlenode{a}{\makebox[0.2cm][c]{\small $E$}}&\\[0.7cm]
&\circlenode{b}{\makebox[0.2cm][c]{\small $M$}}&&\rnode{c}{\small$2$}\\
&&&{\small$0$}\\[0.4cm]
\rnode{d}{\small$-1$}&&\rnode{e}{\small$1$}&\\
{\small$-1$}&&{\small$1$}&\\[0.7cm]
\ncline[linewidth=1pt,linecolor=yellow]{-}{a}{b}\bput[1pt]{0}(0.5){\small$E$}%
\ncline[linewidth=1pt,linecolor=yellow]{-}{a}{c}\aput[1pt]{0}(0.5){\small$N$}%
\ncline[linewidth=1pt,linecolor=yellow]{-}{b}{d}\bput[1pt]{0}(0.5){\small$C$}%
\ncline[linewidth=1pt,linecolor=yellow]{-}{b}{e}\aput[1pt]{0}(0.5){\small$A$}%
%\rput{45}{\psellipse[linecolor=green](4.2,1)(1.3,0.3)}
\end{tabular}
\end{center}

\section{Jeux à information parfaite et complète}

\begin{definition}
C'est un jeu sous forme extensive dans lequel chaque ensemble
d'information est un singleton.
\end{definition}
Des exemples sont fournis par jeu de Nim, d'échecs, de Go, et par
le jeu d'entrée précédent.

Revenons sur ce jeu d'entrée. La menace de jouer $c$ si $E$ joue $e$ est-elle crédible~?\\
\begin{itemize}
\item Supposons que $E$ ait déjà joué $e$. $M$ a-t-il intérêt
à jouer $c$ ou $a$~? ($E$ est ``mis devant le fait accompli'').\\
\item Ayant effectué ce raisonnement, $E$ a-t-il intérêt à jouer
$e$ ou bien $s$~?
\end{itemize}

\subsection{Procédure d'induction à rebours} Comme pour la
programmation dynamique (jeux à un seul joueur). On définit les
stratégies des joueurs en partant des n{\oe}uds terminaux, puis
pour les décisions plus ``en amont'' de l'arbre progressivement
jusqu'au n{\oe}ud initial.

On obtient une stratégie unique pour chaque joueur si pas
d'égalité entre deux paiements de n{\oe}uds terminaux et s'il n'y
a pas de n{\oe}uds où joue la nature.

\begin{theoreme}[dit de Zermelo]
Tout jeu à information complète et parfaite admet un équilibre de
Nash en stratégies pures, qui peut \^etre calculé par la procédure
d'induction à rebours.
\end{theoreme}

 Application au jeu d'échecs : Dans le cas du jeu d'échecs,
appliquer le théorème de Zermelo donne le résultat suivant :
\begin{itemize}
\item Soit un joueur a une stratégie gagnante ;
\item Soit chacun des joueurs peut se garantir la nulle.
\end{itemize}
Pourquoi joue-t-on encore aux échecs ?

\section{\'Equilibre parfait dans les sous-jeux}
La procédure d'induction à rebours dans les jeux à information
complète et parfaite permet de résoudre les jeux en ayant des
prédictions plus fortes que celles de l'équilibre de Nash. Cette
procédure et ces raisonnements peuvent-ils être généralisés aux
jeux où les ensembles d'information ne sont pas des singletons ?

Nous commençons par définir la notion de sous-jeu, qui peut-être
vu comme une partie d'un jeu pouvant être étudiée séparément du
jeu dans son entièreté.

\begin{definition} Un sous-jeu $\Gamma'$ de $\Gamma$ consiste en
\begin{itemize}
\item Un sous-ensemble de n{\oe}uds non terminaux $T'\subset T$ tel que $T' = S(t')$ pour
un certain $t'$ et $\forall t\in T'$, $h(t)\subset T'$.
\item Les joueurs, mouvements de la nature, actions possibles, successeurs,
utilités et ensembles d'information de $\Gamma'$ sont définis
comme dans $\Gamma$.
\end{itemize}
\end{definition}

{\bf Remarque} : En particulier, $\Gamma$ est un sous-jeu de
$\Gamma$. Tout jeu admet par conséquent au moins un sous-jeu :
lui-même. Les sous-jeux de $\Gamma$ non égaux à $\Gamma$ sont
appelés sous-jeux stricts de $\Gamma$.

{\bf Remarque} : Dans un jeu à information parfaite et complète,
on a autant de sous-jeux que de n{\oe}uds non terminaux..

Si on a des stratégies de comportement pour les joueurs dans un
jeu, on connaît leurs probabilités de choix à tous les n{\oe}uds.
Ceci définit en particulier des stratégies dans tous les
sous-jeux.

\begin{definition} Un profil de stratégies $s$ (pures ou de
comportement) est un équilibre parfait dans les sous-jeux de
$\Gamma$ si $s$ induit un équilibre de Nash dans tout les
sous-jeux de $\Gamma$.
\end{definition}

La proposition suivante est immédiate, car le jeu tout entier est
toujours un sous-jeu de lui-même.

\begin{proposition} Un équilibre parfait dans les sous-jeux est
un équilibre de Nash de $\Gamma$. %
\end{proposition}

\begin{proposition}
Tout jeu $\Gamma$ admet un équilibre de Nash parfait dans les
sous-jeux.
\end{proposition}

{\bf Méthode} : Appliquer la procédure d'induction à rebours à
chaque fois que l'on peut.
\begin{enumerate}
\item Identifier un sous-jeu ``minimal'' $\Gamma'$;
\item Calculer un équilibre de Nash de $\Gamma'$ ;
\item Remplacer $\Gamma'$ par ce paiement de Nash ;
\item Recommencer en $1$ tant que possible.
\end{enumerate}
